\documentclass[11pt]{letter}
\usepackage[T1]{fontenc}
\usepackage[margin=1in]{geometry}
\begin{document}
\begin{letter}{Editor\\International Journal of Engineering Science}
\opening{Dear Editor,}

Please consider our manuscript, "Thermopiezoelectric voltage polarity in a BaTiO$_3$ microcantilever: Influence of constrained transverse thermal expansion," for publication in \emph{International Journal of Engineering Science}.

\textbf{Relevance to the journal.} The manuscript addresses a constitutive-reduction issue in the coupled thermopiezoelectric response of an anisotropic solid. It derives the plane-strain thermal loading consistently from three-dimensional constitutive tensors and examines the resulting electrical response using a verified finite-element boundary-value formulation. This combination of continuum mechanics, multiphysics coupling and numerical verification is directly relevant to the \emph{International Journal of Engineering Science}.

\textbf{Contribution of the study.} The main finding is that the plane-strain thermal-stress vector must retain the contribution from constrained transverse thermal expansion. Omitting it changes the predicted piezoelectric--thermal voltage from approximately $+0.19$~V to $-0.22$~V for the BaTiO$_3$ cantilever considered. The study also distinguishes the direct pyroelectric and strain-mediated responses, examines the effect of the root constraint, and quantifies sensitivity to the fitted thermal-expansion coefficients. The contribution lies in identifying and verifying this polarity error in a coupled finite-element calculation, not in proposing a new constitutive law or element.

Thank you for considering our manuscript.

\closing{Sincerely,\\Dr. Anjali Patel\\Corresponding author\\p\_anjali@blr.amrita.edu}
\end{letter}
\end{document}
